\begin{table}[htbp]
   
   \caption{\label{tab:6yr_huc02fe_inorg_ravalli_2005_k2nmines} Effect of the number of upstream coal mines on inorganic chemicals by upstream distance, SYR2 1998--2005}
   \bigskip
   
   \centering
   
\begin{adjustbox}{width = \textwidth, center}
   \begin{tabular}{lcccc}
      \toprule
       & Arsenic & Nitrate & Barium & Selenium \\
                                & (1)      & (2)      & (3)          & (4)\\  
      \midrule 
      \multicolumn{5}{l}{\textbf{Panel A: Within one HUC12 upstream}} \\
      Upstream coal mines (sum) & 0.0000   & 0.0472   & 0.0035$^{*}$ & 0.0001\\   
                                & (0.0001) & (0.0495) & (0.0018)     & (0.0002)\\   
       \\
      \midrule 
      \multicolumn{5}{l}{\textbf{Panel B: Within two HUC12 upstream}} \\
      Upstream coal mines (sum) & 0.0001   & 0.0263   & 0.0019   & -0.0003\\   
                                & (0.0001) & (0.0319) & (0.0018) & (0.0004)\\   
       \\
      Observations              & 366      & 472      & 353      & 350\\  
      \bottomrule
   \end{tabular}
\end{adjustbox}
   
   {\tiny\linespread{1}\selectfont \par \raggedright 
   \textit{Notes:} Within each chemical, the column shows mean measured concentration from the EPA 6-Year Review. Non-detect values replaced by MDL$/\sqrt{2}$ following Ravalli et al.~(2022). Explanatory variable is the number of coal mines in the year in watersheds within one flow step upstream of the utility's intake (Panel A) or within two flow steps upstream (Panel B), summed across those watersheds. Sample: utilities with a coal mine within two flow steps upstream of their intake and no coal mine colocated with their intake. Standard errors clustered at the utility level. All specifications include utility and HUC02 $\times$ year fixed effects. *** p$<$0.01, ** p$<$0.05, * p$<$0.1.}
   
\end{table}


